% mazzi: 06
\chapter{Il tessuto connettivo specializzato}

% ---------------------------------------------------------------------------
\section{Il tessuto cartilagineo}

\subsection{Generalità}
\textbf{Cartilagine}: materiale filogeneticamente antico (nei pesci cartilaginei, es.\ squali,
forma l'intero scheletro). Nell'uomo persiste dove serve una struttura solida ma deformabile e
resistente alla compressione: setto nasale, laringe e apparato respiratorio, dischi intervertebrali,
padiglione auricolare, articolazioni di spalla/gomito/testa omerale, menischi, sinfisi pubica,
articolazioni di mano e piede.
\begin{itemize}
  \item priva di vasi $\rightarrow$ nutrimento per diffusione dal \textbf{pericondrio} (connettivo
    denso di rivestimento);
  \item \textbf{rigenerazione} molto lenta e spesso incompleta.
\end{itemize}

\subsection{I condrociti}
\textbf{Condrociti}: cellule della cartilagine, forma \textbf{globosa}, alloggiate in cavità della
matrice dette \textit{lacune}. Traggono nutrimento solo per diffusione $\rightarrow$ scarsa
capacità proliferativa e rigenerativa. Tre tipi di cartilagine per composizione della matrice:
\textbf{ialina}, \textbf{fibrosa}, \textbf{elastica}.

\subsection{Cartilagine ialina}
\begin{itemize}
  \item la più diffusa: superfici articolari delle ossa lunghe, cartilagini costali, scheletro
    laringeo, anelli di trachea e bronchi;
  \item costituisce il \textbf{modello cartilagineo} da cui, in fase prenatale, si sviluppa gran
    parte dello scheletro per \textbf{ossificazione endocondrale};
  \item matrice omogenea al microscopio ottico, ricca di collagene di tipo II e proteoglicani
    $\rightarrow$ resistenza alla compressione e basso attrito articolare.
\end{itemize}

\subsection{Cartilagine fibrosa}
\begin{itemize}
  \item forma di transizione tra connettivo denso e cartilagine: grossi fasci fibrosi in scarsa
    matrice (con proteoglicani variabili);
  \item nei dischi interarticolari, in particolare i \textbf{menischi} (ben fissata all'osso, più
    resistente alla trazione della ialina).
\end{itemize}

\subsection{Cartilagine elastica}
\begin{itemize}
  \item molto rara: padiglione auricolare e tuba uditiva (di Eustachio);
  \item numerose fibre elastiche nella matrice $\rightarrow$ \textbf{deformabilità} (resiste a urti
    improvvisi e si flette senza rompersi).
\end{itemize}

% ---------------------------------------------------------------------------
\section{Il tessuto osseo}

\subsection{Generalità}
\textbf{Osso}: parte anatomica solida dei vertebrati; l'insieme = \textbf{scheletro}. Ottima
resistenza a compressione e taglio con peso modesto $\rightarrow$ materiale ideale per lo scheletro.
\begin{itemize}
  \item alla nascita $\sim$270 ossa $\rightarrow$ nell'adulto \textbf{206} (68 articolazioni), per
    fusione di alcune ossa durante lo sviluppo;
  \item peso $\sim$20\% del corporeo;
  \item connettivo specializzato con matrice calcificata che imprigiona le cellule; tessuto
    dinamico, si rimodella secondo le forze applicate;
  \item rivestito dal \textbf{periostio} (esterno) e delimitato dall'\textbf{endostio} (interno);
    cavità centrale col midollo (emopoiesi).
\end{itemize}

\subsection{Funzioni del tessuto osseo}
\rubrica{Meccaniche}
\begin{itemize}
  \item \textbf{protettiva}: organi interni (cranio $\rightarrow$ encefalo; gabbia toracica
    $\rightarrow$ cuore e polmoni);
  \item \textbf{motoria}: leve per il movimento;
  \item \textbf{sensoriale}: ossicini dell'udito trasmettono il suono;
  \item \textbf{sostenitiva}: impalcatura del corpo, attacco ai tessuti molli.
\end{itemize}
\rubrica{Emopoietiche}
Midollo delle ossa lunghe e spugnoso delle piatte $\rightarrow$ cellule staminali emopoietiche
(eritrociti e leucociti). Tre tipi di midollo:
\begin{itemize}
  \item \textbf{rosso}: ossa del neonato e, nell'adulto, spugnoso di ossa brevi/piatte ed epifisi
    di femore e omero $\rightarrow$ sede attiva dell'emopoiesi;
  \item \textbf{giallo}: diafisi delle ossa lunghe dell'adulto, ricco di adipociti (riserva
    energetica);
  \item \textbf{grigio}: nell'anziano, inattivo.
\end{itemize}
\rubrica{Metaboliche}
\begin{itemize}
  \item \textbf{riserva di minerali}: calcio e fosforo;
  \item \textbf{riserva di grassi}: acidi grassi nel midollo giallo;
  \item \textbf{riserva di fattori di crescita}: nella matrice (es.\ proteina morfogenetica ossea),
    rilasciati in caso di frattura;
  \item \textbf{detossificazione}: la componente inorganica assorbe metalli pesanti;
  \item \textbf{equilibrio acido-base}: tampona il pH assorbendo/rilasciando sali e ioni;
  \item \textbf{secrezione endocrina}: controllo del metabolismo del fosforo.
\end{itemize}

\subsection{Osso compatto e osso spugnoso}
\textbf{Osso compatto}: porzione esterna delle ossa lunghe; solido, con spazi solo per cellule,
processi e vasi. Formato da \textbf{osteoni} (colonne parallele all'asse) fatti di \textbf{lamelle}
concentriche (\textit{osso lamellare}) attorno a un canale centrale.
\begin{itemize}
  \item \textbf{canale di Havers}: centrale, con vasi e nervi;
  \item \textbf{canale di Volkmann}: trasversale, collega i canali di Havers adiacenti;
  \item \textbf{lamelle}: cerchi concentrici (deposizione apposizionale);
  \item \textbf{lacuna}: ospita l'osteocita;
  \item \textbf{canalicoli}: accolgono i processi degli osteociti $\rightarrow$ scambio di
    metaboliti tra cellule.
\end{itemize}
\textbf{Osso spugnoso}: epifisi e interno delle ossa piatte/brevi; rete di \textbf{trabecole}
sottili con spazi per il midollo. Niente sistemi di Havers $\rightarrow$ scambio coi sinusoidi
midollari tramite i canalicoli. Trabecole rivestite dall'\textbf{endostio} (cellule
osteoprogenitrici, osteoblasti, osteoclasti). Al ME a scansione: compatto denso (canali di Havers)
vs spugnoso reticolare e poroso.

\subsection{Periostio ed endostio}
\begin{itemize}
  \item \textbf{periostio}: connettivo denso irregolare ricco di collagene, circonda l'osso; si
    inserisce con le \textbf{fibre di Sharpey}. Due strati:
    \begin{itemize}
      \item \textit{fibroso esterno}: veicola vasi e nervi;
      \item \textit{cellulare interno}: cellule osteoprogenitrici $\rightarrow$ osteoblasti;
    \end{itemize}
  \item \textbf{endostio}: tappezza le cavità interne (midollare, trabecole, canali di Havers e
    Volkmann); unico strato di cellule osteoprogenitrici, più sottile del periostio; nutrimento e
    produzione di nuove cellule ossee.
\end{itemize}

\subsection{Ossa piatte}
Volta cranica: \textbf{tavolato interno} + \textbf{tavolato esterno} di osso compatto che
racchiudono uno strato di spugnoso (\textbf{diploe}). Esterno rivestito dal \textbf{pericranio},
interno dalla \textbf{dura madre} (che riveste il periostio e protegge l'encefalo).

\subsection{Osteogenesi}
Due modalità di ossificazione prenatale:
\begin{itemize}
  \item \textbf{endocondrale}: indiretta, attraverso un \textbf{modello cartilagineo} (gran parte
    dello scheletro);
  \item \textbf{intramembranosa}: diretta da abbozzo connettivale, senza fase cartilaginea (ossa
    piatte della volta cranica).
\end{itemize}

\rubrica{Ossificazione endocondrale}
\begin{itemize}
  \item abbozzo cartilagineo che cresce per via interstiziale e apposizionale;
  \item dalla 6\textsuperscript{a}--9\textsuperscript{a} settimana fetale: la cartilagine è demolita
    e sostituita da un manicotto osseo, a partire dal \textbf{centro primario} (diafisi);
  \item i condrociti centrali si ipertrofizzano, accumulano glicogeno, si vacuolizzano $\rightarrow$
    lacune più grandi, matrice ridotta e calcificante $\rightarrow$ \textbf{complessi calcificati
    cartilagine/osso} in riassorbimento;
  \item poi compaiono i \textbf{centri secondari} nelle epifisi (senza manicotto osseo): cellule
    osteoprogenitrici $\rightarrow$ osteoblasti $\rightarrow$ matrice ossea; non coinvolgono le
    superfici articolari né la piastra epifisaria (cartilagine di accrescimento);
  \item alla pubertà, per gli ormoni sessuali, le commissure epifisarie si ossificano $\rightarrow$
    fine della crescita in lunghezza e picco di massa ossea (poi ridotta con l'età, più nella donna
    dopo la menopausa).
\end{itemize}

\rubrica{Osteogenesi e vitamina D}
\textbf{Vitamina D}: fondamentale (influenza il bilancio del calcio). Carenza possibile nei
lungodegenti (scarsa esposizione UV per la sintesi cutanea); usata nella profilassi
dell'osteoporosi e nell'alimentazione infantile (con calcio). Pochi alimenti la contengono: olio di
fegato di merluzzo, pesci grassi (salmone, aringa), latte e derivati, uova, fegato, verdure a
foglia verde.

\rubrica{Osteogenesi imperfetta}
Malattia genetica rara (M/F = 1:1) $\rightarrow$ \textbf{fragilità ossea}. 7 varianti, 4 tipi
principali:
\begin{itemize}
  \item \textbf{tipo I}: ritardo di accrescimento, fratture, cifosi, scoliosi, perdita dell'udito;
  \item \textbf{tipo II}: fatale in utero, fratture multiple prenatali;
  \item \textbf{tipo III}: fratture alla nascita, deformazioni progressive, bassa statura;
  \item \textbf{tipo IV}: meno grave, statura poco ridotta, fragilità lieve/moderata.
\end{itemize}

\subsection{Le cellule dell'osso}
Quattro tipi: \textbf{osteoprogenitrici}, \textbf{osteoblasti}, \textbf{osteociti},
\textbf{osteoclasti}.

\rubrica{Cellule osteoprogenitrici}
Dalle cellule mesenchimali embrionali; mantengono la capacità di dividersi; nello strato interno
del periostio e nell'endostio $\rightarrow$ per differenziamento originano gli osteoblasti.

\rubrica{Osteoblasti}
\begin{itemize}
  \item sintetizzano e secernono i costituenti organici della matrice;
  \item cellule cubiche/cilindriche in superficie, nucleo eccentrico, RER e Golgi molto sviluppati;
    giunzioni comunicanti con le cellule vicine;
  \item quando cessano di produrre matrice $\rightarrow$ \textbf{cellule delimitanti l'osso}
    (quiescenti, incapaci di dividersi, ma possono riprendere la secrezione).
\end{itemize}

\rubrica{Osteociti}
\begin{itemize}
  \item cellule mature dell'osso (20.000--30.000/mm\textsuperscript{3}), da osteoblasti intrappolati
    nelle lacune;
  \item più piccoli e meno basofili degli osteoblasti, nucleo appiattito, pochi organuli;
  \item processi nei canalicoli, uniti da giunzioni comunicanti (gap-junction);
  \item mantengono la matrice, rilasciano $Ca$ quando serve, facilitano i preosteoblasti nel
    rimodellamento.
\end{itemize}

\rubrica{Osteoclasti}
\begin{itemize}
  \item cellule multinucleate da progenitori macrofagico-monocitari del midollo;
  \item riassorbono la matrice sotto \textbf{ormone paratiroideo} e \textbf{calcitonina}
    (decalcificazione, digestione, assorbimento);
  \item occupano le \textbf{lacune di Howship} (depressioni nella superficie in riassorbimento).
\end{itemize}

\subsection{Equilibrio tra osteoblasti e osteoclasti}
Riassorbimento e neoformazione devono bilanciarsi: l'osso distrutto dagli osteoclasti va rimpiazzato
da quello prodotto dagli osteoblasti. Processo regolato da cellule, alterabile con l'invecchiamento
o lo sbilanciamento dei controlli. L'osso è al servizio delle necessità di calcio (prioritario per
nervo e muscolo) $\rightarrow$ il carico meccanico quotidiano è essenziale per mantenerne massa e
struttura.
